\section{Method}
\label{sec:method}

\subsection{Problem Definition}
\label{sec:method:problem}

Let $f_\theta$ denote a language model and $\mathcal{D}_0$ a corpus of human-written seed text.
We define a \emph{recursive generation chain} of depth $K$ as the sequence $\mathcal{D}_0, \mathcal{D}_1, \dots, \mathcal{D}_K$, where each $\mathcal{D}_k$ ($k \geq 1$) is produced by finetuning $f_\theta$ on $\mathcal{D}_{k-1}$ and sampling from the resulting model.
Depth $d{=}0$ denotes human-written text; depth $d{=}k$ ($k \geq 1$) denotes text generated by a model trained on depth-$(k{-}1)$ data.

Given a text passage $x$, the task is to estimate its recursive generation depth $d \in \{0, 1, \dots, K\}$.
We use \emph{depth estimation} as shorthand throughout; the underlying task is ordinal classification.
Depths carry a natural ordering, and $d{=}1$ is structurally closer to $d{=}0$ than $d{=}3$ is, since each generation step compounds distributional drift.

\subsection{Multi-Generational Benchmark Construction}
\label{sec:method:benchmark}

We construct a benchmark spanning three axes: model family, decoding strategy, and generation depth.

\paragraph{Model families.}
We select three 1--1.5B-parameter models with diverse pretraining corpora and architectures: Pythia-1.4B (EleutherAI), OLMo-1B (AI2), and GPT-2~XL (OpenAI, 1.5B).
Restricting to a narrow parameter range isolates the effect of architecture and training data from scale.

\paragraph{Decoding strategies.}
Each model is paired with three decoding strategies: nucleus sampling ($p{=}0.95$) \cite{holtzman2020_nucleus}, temperature sampling ($\tau{=}0.9$), and top-$k$ sampling ($k{=}50$).
These span the dominant modes of open-ended generation and induce qualitatively different entropy profiles.
Our findings on feature importance and discriminability are limited to this grid; transfer to untested models, scales, or strategies is not established.

\paragraph{Depth and scale.}
The primary benchmark covers depths d0--d3 across all $3 \times 3 = 9$ model--strategy cells, with 5{,}000 samples per depth per cell (180K cell-instances; 140K unique passages, as d0 samples are drawn from a shared Wikipedia seed corpus across all cells).
For depth extension experiments (\S\ref{sec:atlas}), we extend all three models under nucleus sampling to d0--d5.

\paragraph{Generation pipeline.}
Starting from a Wikipedia subset ($\mathcal{D}_0$), each cycle: (1)~LoRA finetuning \cite{hu2022lora} on $\mathcal{D}_{k-1}$ ($r{=}16$, $\alpha{=}32$, lr $2{\times}10^{-4}$, effective batch 16, 3 epochs, dropout $0.05$; fresh adapter from the base model at each depth; full protocol in Appendix~\ref{app:lora_protocol}); (2)~sampling 5{,}000 passages (max 512 tokens, beginning-of-sequence (BOS) seeded) under the specified strategy; (3)~$\mathcal{D}_k$ feeds the next cycle.
Each depth is produced by a model trained on prior-depth data, not by sampling from the unmodified base model.
BOS seeding may amplify depth signal; prompted-generation validation (Appendix~\ref{app:prompted_gen}) confirms signal persistence albeit attenuated ($-$13.8\,pp).
No per-passage lineage exists across depths; cross-validation splits do not introduce leakage.

\paragraph{Reference model.}
All statistical features (\S\ref{sec:method:features}) use Pythia-1.4B as fixed reference model.
This overlaps with one generator (Pythia), conferring a 4--11\,pp self-reference advantage on that pipeline (Appendix~\ref{app:ref_ablation}); Pythia results in Table~\ref{tab:main_matrix} are marked $^\dagger$ accordingly.
Alternative references yield 67--79\%, all exceeding $2.5{\times}$ chance (Limitations).

\subsection{Statistical Feature Design}
\label{sec:method:features}

We extract a feature vector $\mathbf{f}(x)$ comprising 19 candidate features from the reference-model perplexity distribution over each passage $x$; ablation (\S\ref{sec:atlas:features}) shows the top 5 achieve 97.6\% of full-set performance, and variance inflation analysis reveals VIF${>}$10 for 14--16 features; effective independent dimensionality is ${\approx}$4--5.
The features fall into four categories (Table~\ref{tab:features}).

\begin{table}[t]
\centering
\small
\caption{The 19 statistical features comprising $\mathbf{f}(x)$, grouped by category. All perplexity and surprisal features are computed from the reference model (Pythia-1.4B).}
\label{tab:features}
\begin{tabular*}{\columnwidth}{@{\extracolsep{\fill}}llr}
\toprule
Category & Features & $n$ \\
\midrule
Perplexity & mean, var, skewness, kurtosis, & 9 \\
statistics & p10, p25, p50, p75, p90 & \\
\addlinespace
Surprisal & mean, var, entropy & 3 \\
statistics & & \\
\addlinespace
Lexical & type--token ratio, hapax ratio, & 4 \\
diversity & bigram entropy, trigram entropy & \\
\addlinespace
Repetition & repeated 2/3/4-gram fractions & 3 \\
\bottomrule
\end{tabular*}
\end{table}

\emph{Perplexity statistics} (9 features) capture the per-token perplexity distribution via moments and quantiles (p10--p90); tail quantiles are especially informative as recursive finetuning erodes the high-perplexity tail (\S\ref{sec:atlas}).
\emph{Surprisal statistics} (3 features): mean, variance, and entropy of $-\log p(\text{token})$.
\emph{Lexical diversity} (4 features): type--token ratio, hapax ratio, bigram/trigram entropy.
\emph{Repetition} (3 features): repeated $n$-gram fractions for $n \in \{2, 3, 4\}$.

\subsection{Effective Discriminable Depth}
\label{sec:method:edd}

Aggregate multi-class accuracy obscures whether discriminability degrades at specific depth boundaries.
We propose the \emph{Effective Discriminable Depth} (EDD) as an operational measure of how far into a recursive chain adjacent depths remain distinguishable.

\paragraph{Pairwise discriminability.}
Let $\mathcal{D} = \{0, 1, \dots, K\}$ denote the depth set and $T \in (0.5, 1.0]$ a discrimination threshold.
For each adjacent pair $(d, d{+}1)$, we define the pairwise discriminability $\mathrm{AUC}(d, d{+}1)$ as the area under the ROC curve of a dedicated binary Random Forest classifier (500 trees, $\lfloor\sqrt{19}\rfloor$ candidate features per split) trained on class-balanced data from $\mathbf{f}(x)$ (\S\ref{sec:method:features}) to distinguish depth-$d$ from depth-$(d{+}1)$ samples, evaluated via 5-fold stratified cross-validation.
Each pair uses an independent classifier to avoid decision boundary dilution from multi-class settings.

\begin{definition}[Effective Discriminable Depth]
\label{def:edd}
The \emph{Effective Discriminable Depth} at threshold $T$ is:
\begin{multline}
\mathrm{EDD}_T = \max \bigl\{ d^*{+}1 \,:\, \\
\forall\, d' \leq d^*,\; \mathrm{AUC}(d', d'{+}1) \geq T \bigr\}
\label{eq:edd}
\end{multline}
where $d^* \in \{0, \dots, K{-}1\}$. If\, $\mathrm{AUC}(0,1) < T$, then $\mathrm{EDD}_T = 0$.
\end{definition}

$\mathrm{EDD}_T$ reports the longest \emph{contiguous} discriminable prefix: the deepest $d$ such that every pair from $(0,1)$ through $(d{-}1, d)$ exceeds $T$.
All AUC profiles in the primary benchmark decay monotonically; non-monotonic profiles occur only in cross-scale extensions (Appendix~\ref{app:cross_scale_full}).
We adopt $T{=}0.70$ following the ``acceptable discrimination'' threshold in ROC analysis \cite{hosmer2000} (8/9 cells achieve $\mathrm{EDD}_{0.70} \geq 3$); sensitivity analysis (\S\ref{sec:atlas:edd}) shows $\mathrm{EDD}$ shifts by at most one level across $T \in \{0.65, 0.70, 0.80\}$ in all nine cells.

\paragraph{Classifiers.}
Our primary classifier is a Random Forest (200 trees, unconstrained depth, 5-fold CV; EDD binary classifiers use 500 trees).\footnote{Statistical tests are not corrected for multiple comparisons; $p$-values should be interpreted as exploratory.}
The RF produces ordinally consistent predictions: 98.2\% of errors fall on adjacent depths ($|\Delta| \leq 1$), quadratic weighted kappa (QWK)\,=\,0.89.
Comparisons against other classifiers and detection baselines appear in \S\ref{sec:generalization}.